\section{模型检验与灵敏度分析}

\subsection{统一约束审计与双运行检验}
任务审计逐项检查唯一执行、实时即时开工、到达、截止、单向网络时延、第 2406 小时禁占用、GPU、IT 和设施功率。能源审计统一检查
\begin{align}
e^{energy}_{rt}&=G^{buy}_{rt}+R^{av}_{rt}+D_{rt}-L^F_{rt}-C_{rt}-G^{sell}_{rt}-R^{curt}_{rt},\\
e^{SOC}_{rt}&=S_{rt}-S_{r,t-1}-\eta_r^cC_{rt}+D_{rt}/\eta_r^d.
\end{align}
此外逐行复核五个来源分解恒等式、$N_{rt}=G^{buy}_{rt}-G^{sell}_{rt}$、购售电与充放电互斥、功率边界和终端 SOC 等式。solve v19 中任务容量审计的最大正残差为 $1.364\times10^{-12}$；Q2 能源重算最大绝对残差为 $1.695\times10^{-11}$；Q3 与 Q4 的能源平衡和来源分解最大绝对残差分别为 $5.684\times10^{-14}$ MW 与 $1.421\times10^{-14}$ MW，其余 SOC、互斥、净购电恒等式和终端残差均为 0。

五方向 6 遍共形成 30 个累计状态检查点。候选级确定性分层审计共 430 个样本，覆盖 6 个 pass、accepted、feasible-but-rejected、capacity-rejected、no-change、fractional-duration、跨 2399/2400、返回基线、高功率与高 GPU 等类别；最大绝对和相对残差分别为 $8.009\times10^{-8}$ 和 $2.208\times10^{-9}$。公共候选池的 CandidateID 与语义三元组均唯一，全部 50000 个任务被检查，具有有效替代候选的任务覆盖率为 99.024\%。

两次完整正式复算已在本次论文修订之前由 code v21 完成，自然耗时分别为 4558.731 s 和 3786.555 s，均未设置显式墙钟上限。两次的 code SHA-256 均为 \texttt{1c777886...a6b71}，\texttt{results.json} SHA-256 均为 \texttt{c63996f2...9389e}；比较 106 份文件后结构化语义差异为 0。本次 paper-only 修订不再次运行 Q1--Q4 求解。

\subsection{问题二完整邻域多遍稳定性}
120000 行基础候选池的 BaseCandidateID 与由 TaskID、Region、CandidateStart 组成的语义三元组均 100\% 唯一。每个方向每遍按相同顺序完成 120000 次尝试；成本、直接碳、新能源、服务和平衡分别在第 4、2、4、2、5 遍达到预声明的主目标改善与接受率阈值，故正式共同完成点为第 5 遍。第 6 遍成本和平衡仍分别接受 4 和 1 个极小改善，说明此处“converged”是有限候选邻域的阈值稳定，不是零接受固定点或全局最优。

\subsection{统一参数灵敏度}
平衡方案基准成本为 $-4.393037\times10^8$ 元，直接购电量为 1.5264 MWh。购电价格整体缩放到 0.8、0.9、1.1、1.2 时，净成本分别为 $-4.393038$、$-4.393037$、$-4.393037$、$-4.393036$ 亿元；相对基准只变化约 $-79.37$ 至 $79.37$ 元，反映当前接近零购电边界，而非一般意义上的价格不敏感。新能源水平系数为 0.8、0.9、1.1、1.2 时，成本分别为 $-3.169013$、$-3.998509$、$-4.544149$、$-4.588926$ 亿元；新能源下降时购电缺口扩大，上升时外送上限与弃电约束削弱边际收益。

图~\ref{fig:sensitivity}a 采用“相对基准成本变化/元”避免用亿元坐标掩盖价格扰动的微小响应；图~\ref{fig:sensitivity}b 保留新能源扰动下的原始净成本量纲。两幅图只复算正式灵敏度 CSV 中已给出的目标值，不重新优化调度；它们与图~\ref{fig:q4scenarios}的联合情景分析口径不同，前者检验 Q2 固定方案参数响应，后者检验 Q4 联合任务--储能情景。
\begin{figure}[H]\centering
\includegraphics[width=.94\textwidth]{fig06_sensitivity.png}
\caption{统一参数灵敏度：（a）购电价格变化下相对基准的成本变化；（b）新能源水平变化下的净成本。}\label{fig:sensitivity}
\end{figure}

\subsection{SOC 网格与波动情景}
15、21、31 个 SOC 状态下，所选平衡策略成本分别为 $-4.528843$、$-4.541740$、$-4.559475$ 亿元；21 相对 31 状态的成本与新能源利用率相对误差分别为 0.38898\% 和 0.07163\%，三组最大约束残差均不超过 $5.684\times10^{-14}$。对应等效完整循环为 84.604、92.401、102.615。题面未给退化成本和循环上限，故只作为理想设备网格敏感性。

Q4 三个新能源波动种子的峰值购电为 241.896 至 320.783 MW、碳排放为 1351.780 至 1716.071 tCO$_2$，显示单一随机轨迹不足以概括风险；当前三种子只提供方向性稳健性证据，不能提升 scenario-feasible 可信等级。

